% Part: lambda-calculus
% Chapter: lambda-definability
% Section: fixpoints

\documentclass[../../../include/open-logic-section]{subfiles}

\begin{document}

\olfileid[hi]{lam}{ldf}{fp}

\olsection{अचल-बिंदु}

मान लें हम क्रमगुणित फलन को पुनरावर्तन द्वारा ऐसे पद $\fn{Fac}$ के रूप में
परिभाषित करना चाहें जिसमें निम्न गुण हो:
\[
\fn{Fac} \ident \lambd[n][\fn{IsZero}\, n\, \num 1 (\fn{Mult}\, n (\fn{Fac}(\fn{Pred} \, n)))]
\]
अर्थात् $n = 0$ होने पर $n$ का क्रमगुणित $1$ है, और अन्यथा $n-1$ के
क्रमगुणित का $n$ गुना। निश्चय ही हम पद $\fn{Fac}$ को इस प्रकार परिभाषित
नहीं कर सकते, क्योंकि दाएँ पक्ष में स्वयं $\fn{Fac}$ आता है। स्वसंदर्भ
वाली ऐसी पुनरावर्ती परिभाषाएँ लैम्ब्डा कलन का भाग नहीं हैं। उदाहरणार्थ,
किसी पद को
\[
\fn{Mult} \ident \lambd[ab][a (\fn{Add}\, a) 0]
\]
द्वारा परिभाषित करने में दाएँ पक्ष पर केवल $\fn{Add}$ जैसे पहले से
परिभाषित पद आते हैं। $\fn{Add}$ को उसके परिभाषक पद से बदलकर हम उसे सदा
हटा सकते हैं। इससे पद $\fn{Mult}$ शुद्ध लैम्ब्डा पद के रूप में मिलेगा; यदि
$\fn{Add}$ में स्वयं परिभाषित पद आते हों (जैसे $\fn{Add}'$ में), तो हम इस
प्रक्रिया को जारी रखकर अंततः शुद्ध लैम्ब्डा पद तक पहुँच सकते हैं।

किंतु ऊपर की $\fn{Fac}$ जैसी पुनरावर्ती परिभाषाओं के लिए यह सत्य नहीं है।
यदि दाएँ पक्ष पर $\fn{Fac}$ के आविर्भाव को स्वयं $\fn{Fac}$ की परिभाषा से
बदलें, तो मिलता है:
\begin{align*}
  \fn{Fac} & \ident \lambd[n][\fn{IsZero}\, n\, \num{1}] \\
    & \qquad (\fn{Mult}\, n
  ((\lambd[n][\fn{IsZero} \, n \, \num 1\, (\fn{Mult}\, n\, (\fn{Fac}
    (\fn{Pred}\, n)))]) (\fn{Pred}\, n)))
\end{align*}
और हमने अब भी दाएँ पक्ष से $\fn{Fac}$ नहीं हटाया है। स्पष्टतः, इस प्रक्रिया
को दोहराने पर परिभाषा लंबी होती जाती है और कभी शुद्ध लैम्ब्डा पद नहीं
मिलता। अतः क्रमगुणित (या सामान्यतः पुनरावर्ती फलनों) को परिभाषित करने की
यह विधि व्यावहारिक नहीं है।

फिर भी पुनरावर्ती परिभाषा हमें कुछ बताती है: यदि $f$ क्रमगुणित फलन को
निरूपित करने वाला पद होता, तो पद
\[
\fn{Fac}' \ident \lambd[g][\lambd[n][\fn{IsZero} \, n \, \num 1\, (\fn{Mult} \, n\, (g (\fn{Pred} n)))]]
\]
पद $f$ पर लागू होने पर, अर्थात् $\fn{Fac}'\,f$, क्रमगुणित फलन को भी
निरूपित करता है। अर्थात् यदि $\fn{Fac}'$ को ऐसा फलन मानें जो कोई फलन
स्वीकार करके फलन लौटाता है, तो $f$ के क्रमगुणित होने की दशा में
$\fn{Fac'}\, f$ का मान केवल~$f$ है। जिस फलन~$f$ में $\fn{Fac}' \, f
\equal[\beta] f$ गुण हो, उसे $\fn{Fac}'$ का \emph{अचल-बिंदु} कहते हैं। अतः
क्रमगुणित $\fn{Fac}'$ का अचल-बिंदु है।

लैम्ब्डा कलन में ऐसे पद हैं जो किसी दिए पद के अचल-बिंदुओं की संगणना करते
हैं; फिर इन पदों का उपयोग $\fn{Fac}'$ जैसे पद को क्रमगुणित की परिभाषा में
बदलने के लिए किया जा सकता है।

\begin{defn}
  \ollabel{defn:Turing-Y}
  \emph{Y-संयोजक} यह पद है:
  \[
  Y \ident (\lambd[ux][x(uux)])(\lambd[ux][x(uux)]).
  \]
\end{defn}

\begin{thm}
  $Y$ में यह गुण है कि किसी भी पद $g$ के लिए $Yg \red g(Yg)$। अतः $Yg$
  सदा~$g$ का अचल-बिंदु है।
\end{thm}

\begin{proof}
  $(\lambd[ux][x(uux)])$ को~$U$ से संक्षिप्त करें, ताकि $Y \ident UU$। तब
  \begin{align*}
    Y g &\ident (\lambd[ux][x(uux)])U\,g \\
    &\red (\lambd[x][x(UUx)])g\\
    & \red g(UUg) \ident g(Yg).
  \end{align*}
  चूँकि $g(Yg)$ और $Yg$ दोनों अपचयित होकर $g(Yg)$ बनते हैं, इसलिए $g(Yg)
  \equal[\beta] Yg$; अतः $Yg$, ~$g$ का अचल-बिंदु है।
\end{proof}

निश्चय ही, चूँकि $Yg$ रिडेक्स है, इसलिए अपचयन अनिश्चित काल तक चल सकता है:
\begin{align*}
  Y g &\red g (Y g) \\
    &\red g (g (Y g)) \\
    &\red g(g (g (Y g)))\\
    &\ldots
\end{align*}
अतः हम $Yg$ को स्वयं पर अनंत बार लागू $g$ मान सकते हैं। यदि उस पर~$g$ एक
बार और लागू करें, तो कहने को हम कुछ अतिरिक्त नहीं कर रहे; ~$Yg$ पर अनंत
बार लागू~$g$ पर फिर लागू~$g$, अब भी~$Yg$ पर अनंत बार लागू~$g$ ही है।

ध्यान दें कि~$Yg$ से आरंभ होने वाला ऊपर का $\beta$-अपचयन चरणों का अनुक्रम
अनंत है। अतः यदि हम $Yg$ को किसी पद पर लागू करें, अर्थात् $(Yg)N$ पर विचार
करें, तो वह पद भी अपचयित होकर अनंत रूप से अनेक भिन्न पद बनाएगा:
$(g(Yg))N$, $(g(g(Yg)))N$, \dots। फिर भी संभव है कि अपचयन चरणों का कोई
\emph{अन्य} अनुक्रम किसी प्रसामान्य रूप पर समाप्त हो।

उदाहरण के लिए क्रमगुणित लें। $\fn{Fac}$ को $Y\, \fn{Fac}'$ (अर्थात्
$\fn{Fac'}$ के किसी अचल-बिंदु) के रूप में परिभाषित करें। तब:
\begin{align*}
  \fn{Fac}\,\num 3
  &\red Y\, \fn{Fac}' \, \num 3 \\
  &\red \fn{Fac}' (Y \,\fn{Fac}') \, \num 3 \\
  &\ident (\lambd[x][\lambd[n][\fn{IsZero} \, n\, \num 1\,
      (\fn{Mult}\, n\, (x (\fn{Pred}\, n)))]]) \, \fn{Fac} \, \num 3\\
  & \red \fn{IsZero} \, \num 3\, \num 1\,
      (\fn{Mult}\, \num 3\, (\fn{Fac} (\fn{Pred}\, \num 3))) \\
  &\red  \fn{Mult}\, \num 3 \, (\fn{Fac} \, \num 2).
  \intertext{इसी प्रकार,}
  \fn{Fac}\,\num 2
  &\red  \fn{Mult}\, \num 2 \, (\fn{Fac} \, \num 1) \\
  \fn{Fac}\,\num 1
  &\red  \fn{Mult}\, \num 1 \, (\fn{Fac} \, \num 0)
  \intertext{किंतु}
  \fn{Fac}\,\num 0
  &\red \fn{Fac}' (Y \,\fn{Fac}') \, \num 0 \\
  &\ident (\lambd[x][\lambd[n][\fn{IsZero} \, n\, \num 1\,
      (\fn{Mult}\, n\, (x (\fn{Pred}\, n)))]]) \, \fn{Fac} \, \num 0\\
  & \red \fn{IsZero} \, \num 0\, \num 1\,
      (\fn{Mult}\, \num 0\, (\fn{Fac} (\fn{Pred}\, \num 0))).\\
  &\red \num 1.
  \intertext{अतः सब मिलाकर}
  \fn{Fac}\,\num 3
  &\red  \fn{Mult}\, \num 3 \,
  (\fn{Mult} \, \num 2\, (\fn{Mult} \, \num 1\, \num 1)).
\end{align*}

$\fn{Fac'}$ के लिए जो लागू होता है, वह किसी भी पुनरावर्ती परिभाषा के लिए
लागू होता है। मान लें हमारे पास पुनरावर्ती समीकरण है:
\begin{align*}
g\,x_1\dots x_n & \equal[\beta] N
\intertext{जहाँ $N$ में~$g$ और $x_1$, \dots,~$x_n$ आ सकते हैं। तब सदा कोई
ऐसा पद~$G \ident (Y \lambd[g][\lambd[x_1\dots x_n][N]])$ है कि}
G\,x_1 \dots x_n & \equal[\beta] \Subst{N}{G}{g}.
\intertext{क्योंकि अचल-बिंदु प्रमेय से,}
G \ident (Y \lambd[g][\lambd[x_1\dots x_n][N]]) & \red \lambd[g][\lambd[x_1\dots x_n][N]](Y \lambd[g][\lambd[x_1\dots x_n][N]])\\
& \ident (\lambd[g][\lambd[x_1\dots x_n][N]])\,G
\intertext{और फलतः}
G\,x_1 \dots x_n 
& \red (\lambd[g][\lambd[x_1\dots x_n][N]])\,G\,x_1\dots x_n\\
& \red (\lambd[x_1\dots x_n][\Subst{N}{G}{g}])\,x_1\dots x_n\\
  & \red \Subst{N}{G}{g}.
\end{align*}

\olref{defn:Turing-Y} का $Y$ संयोजक एलन ट्यूरिंग का है। अलॉन्ज़ो चर्च ने
एक भिन्न रूप प्रस्तावित किया था, जिसे हम~$Y_C$ कहेंगे:
\[
Y_C \ident \lambd[g][(\lambd[x][g(xx)])(\lambd[x][g(xx)])].
\]
चर्च का संयोजक ट्यूरिंग के संयोजक से कुछ दुर्बल है: $Yg \equal[\beta]
g(Yg)$, किंतु $Yg \bred g(Yg)$ नहीं। $V$ पद $\lambd[x][g(xx)]$ हो, ताकि
$Y_C \ident \lambd[g][VV]$। तब
\begin{align*}
  VV & \ident (\lambd[x][g(xx)])V \red g(VV)
  \text{ और अतः}\\
  Y_C g & \ident (\lambd[g][VV])g \red VV \red g(VV), \text{ किंतु साथ ही}\\
  g(Y_C g) & \ident g((\lambd[g][VV])g) \red g(VV).
\end{align*}
दूसरे शब्दों में, $Y_Cg$ और $g(Y_Cg)$ अपचयित होकर समान पद $g(VV)$ में
पहुँचते हैं; अतः $Y_Cg \equal[\beta] g(Y_Cg)$। अनुप्रयोगों के लिए यह प्रायः
पर्याप्त है।

\end{document}

